\chapter{まとめと今後の展望}

\section{まとめ}
本研究では，個人の性格特性と日々の推し活の状況から，その人に適した推し活方法を提案するシステムOshiFitを構築し，その有用性を検証した．

第1章では，メンタル不調の増加と推し活の広がりを背景に，推し活が便益と負担の両方をもたらすこと，本人は自分の負担の原因に気づきにくいこと，先行研究には便益と負担を同時に扱った定量的な検討や個人を支援する仕組みが不足していることを述べ，3つのリサーチクエスチョンを設定した．

第2章では，推し活・ファン行動，ウェルビーイングとバーンアウト，性格特性と個人化，パーソナルインフォマティクスと行動変容，アバターと将来の自己，予測の不確かさの提示に関する先行研究を概観し，本研究の位置付けを示した．

第3章では，215名のファンを対象とした質問紙調査により，RQ1に答えた．推し活の頻度と没入は幸福度と推し疲れの両方と関連し，便益と負担が同じ行動から生じることが分かった．性格特性は推し活の関わり方と関連したが，行動の効果を調整してはおらず，性格だけでは明確なタイプに分かれなかった．個人の予測の精度は限られていた（交差検証決定係数：幸福度.27，推し疲れ.11）．これらの結果から，一律に減らす助言をしない（D1），性格と行動を組み合わせて個人化する（D2），便益と負担を同時に示す（D3），不確かさを正直に示す（D4）という4つの設計指針を導いた．

第4章では，設計指針に基づいてOshiFitを設計・実装した．OshiFitは67項目の質問紙から6軸の得点を算出して算出式とともに示し，すでに良好な人には今のままでよいことを伝え，それ以外の人には11種類のプランから3案を示す．さらに，プランAとプランBの6週間後の見込みを幅とともに本人に似たアバターで示し，本人が選んだうえで6週間の振り返りを支援する．

第5章では，RQ2とRQ3に答えるための6週間の検証実験の方法と分析計画を述べた．\todo{実験後に記入：主な結果}

第6章では，便益と負担が同じ行動から生じることの意味，性格特性の役割，不確かさを示すことの意味，アバターで2つの未来を示すことについて考察し，本研究の限界を述べた．

\section{今後の展望}
今後の展望として，次の点が挙げられる．

\begin{enumerate}
\item \textbf{対象の拡大}：年齢層や推しのジャンル（アイドル，キャラクター，VTuber，スポーツ選手など）をより幅広く含めた調査を行い，知見の一般化を図る．
\item \textbf{想定改善量の更新と予測の改善}：検証実験で得られた正解データにより想定改善量を更新し（モデルv1.1），利用者が増えるにつれて予測の精度を高める．週ごとの振り返りのデータを用いた，個人内の変化を捉える予測モデルも検討する．
\item \textbf{アバターの効果の検証}：本研究ではアバターの有無による比較を行っていない．今後は，アバターで見込みを示す場合と文章や数値だけで示す場合を比較し，アバターそのものの効果を検証する．
\item \textbf{生理指標との組み合わせ}：前研究で用いたスマートウォッチによる心拍変動\cite{thayer2007}などの生理指標を組み合わせ，自己報告だけでは捉えにくいストレスの変化を捉える．
\item \textbf{長期的な効果}：6週間より長い期間の利用で，プランが習慣として定着するか，推し疲れの低下が持続するかを検証する．
\item \textbf{社会実装}：研究で構築したOshiFitを，誰もが気軽に使えるアプリケーションとして公開し，展示やファンコミュニティとの連携を通じて，推し活を持続可能な形で楽しむための支援として普及させる．
\end{enumerate}
